% Section 2. Sources: 637_rung2_matrix_product/PROOF.md section 1 (Proposition T) and FINAL.md;
% 677_tape_bijection_review/REVIEW.md; 607_bounds_paper/paper_word.tex Appendix A (labels, stacks, words) and
% Appendix B (linear completion); 805_entropy_certificates/PROOF.md section 1.2 step 3 (Lemmas R and B);
% 641_tape_window_bounds/PROOF.md section 0 (class C); web/src/meandric/tapeMachine.ts and TapeMachineWriteup.tsx
% (rules as implemented, gap numbers, picture of a move). Counts: 1251 FINAL (n <= 9), 1252 FINAL EY6 (n <= 15).
% The worked examples were re-derived with a small script written for this paper (rules of tapeMachine.ts).
%%%%%%%%%%%%%%%%%%%%%%%%%%%%%%%%%%%%%%%%%%%%%%%%%%%%%%%%%%%%%%%%%%%%%%%%
\section{The tape machine}\label{sec:tape}
%%%%%%%%%%%%%%%%%%%%%%%%%%%%%%%%%%%%%%%%%%%%%%%%%%%%%%%%%%%%%%%%%%%%%%%%

Number the crossings $1,\ldots,2n$ from left to right. The river cuts the road into $2n$ pieces, the \emph{arches},
each joining two crossings: $n$ \emph{upper} arches above the river and $n$ \emph{lower} arches below it. Arches on the
same side do not cross. In the running example
(Figure~\ref{fig:cut}) the upper arches join $1$--$2$, $3$--$6$, $4$--$5$ and the lower arches join $1$--$6$,
$2$--$5$, $3$--$4$.

\subsection{Labels}

Delete crossing $2n$, that is, cut the road open there. The road becomes a path through the crossings
$1,\ldots,2n-1$, whose two ends are the arches that went to $2n$. Give crossing $1$ the label $g(1)=0$ and continue
along the path with the rule
\begin{equation}\label{eq:arch}
  g(j')-g(j)=\begin{cases}-1&\text{for an upper arch from crossing $j$ to crossing $j'$, $j<j'$},\\ +1&\text{for a lower arch from crossing $j$ to crossing $j'$, $j<j'$}.\end{cases}
\end{equation}
The path has no cycle, so $g$ is well defined. Intuitively, $g(j)$ counts the half turns of the road on its way from
crossing $1$ to crossing $j$: read from left to right, an arch below the river turns one way and an arch above turns
the other. In the running example the labels of crossings $1,\ldots,5$ are $0,-1,0,1,0$.

\begin{lemma}[labels]\label{lem:labels}
For a meander, neighboring crossings $j$ and $j+1$ along the river ($1\le j\le 2n-2$) have labels that differ by
exactly one, and $g(2n-1)=0$. Conversely, suppose that $n$ noncrossing upper arches and $n$ noncrossing lower arches
on the points $1,\ldots,2n$ admit a function $g$ on
$1,\ldots,2n-1$ that obeys \eqref{eq:arch} on every arch not ending at $2n$ and changes by $\pm1$ between neighbors,
then the arches form a single closed road.
\end{lemma}

\begin{proof}[Idea of the proof]
There are $2n-2$ river segments between neighboring crossings $j,j+1\le 2n-1$. Above the river, each of the $n-1$
upper arches not ending at $2n$ bounds a \emph{face}: the region between that arch, the river and the arches directly
under it. Below the river there are $n-1$ such faces as well. For a single road, the road cuts the sphere into two
disks, the river segments in each disk are disjoint chords of it, and the faces in each disk are arranged as a tree.
This lets one assign each river segment to exactly one face: to the face above it if the label steps down along it, to
the face below it if the label steps up. Every face receives exactly one segment. Adding up the label steps along the
bottom of a face, over its river segments and across the arches under it, gives the rule \eqref{eq:arch} for the arch
on top of it. For the converse,
suppose the arches formed a second closed component. It bounds a disk that contains no other component; the river
cuts this disk into one more face than it has segments inside, but each face needs a segment of its own. The full
argument is in the repository (Appendix~\ref{app:where}).
\end{proof}

\subsection{The two stacks at a cut}

Cover the part of the meander to the right of a vertical line, the \emph{cut}, placed between crossings $j$ and $j+1$.
Intuitively, this is like sliding a piece of paper from left to right along the meander, stopping between two
crossings. The arches that start at or before crossing $j$ and end after the cut are \emph{open}. The open arches on
one side are nested, so the innermost one closes first: they form a \emph{stack}, a list in which the last entry added
is the first one removed. Let
$\sigma=(\sigma_1,\ldots,\sigma_a)$ be the labels of the left ends of the open lower arches, from the outermost to
the innermost, and $\tau=(\tau_1,\ldots,\tau_b)$ the same for the open upper arches. The \emph{stack state} is
$(g(j),\sigma,\tau)$. After the first crossing it is $(0,(0),(0))$: one arch leaves crossing~$1$ on each side.

Revealing crossing $j+1$ works as follows. Its label is $g(j)\pm1$ by Lemma~\ref{lem:labels}. On each side of the
river the new crossing either \emph{opens} an arch, whose label $g(j+1)$ is pushed on the stack, or \emph{closes} the
innermost open arch, which is popped. A closing must agree with \eqref{eq:arch}: a lower arch can close only if its
label is $g(j+1)-1$, an upper arch only if its label is $g(j+1)+1$. One more condition comes from
Lemma~\ref{lem:labels}: the label can step up only if the innermost open lower arch has label $g(j)$, and down only if
the innermost open upper arch has label $g(j)$. We call these rules the \emph{label machine}. With the help of the
converse half of Lemma~\ref{lem:labels}, the sequences of updates that the label machine allows correspond exactly to
the meanders (Appendix~\ref{app:where}).

The stack states that occur have four properties: neighboring entries of a stack differ by one; the innermost lower
label is $g(j)$ or $g(j)-2$; the innermost upper label is $g(j)$ or $g(j)+2$; and the two innermost labels are not
both two away from $g(j)$. A stack is empty only in two situations: the lower stack can be empty only when $g(j)=1$,
and the upper stack only when $g(j)=-1$.


\subsection{Tapes}

Write the numbers
\begin{equation}\label{eq:path}
  -1,\ \sigma_1,\ \ldots,\ \sigma_a,\ \tau_b-1,\ \ldots,\ \tau_1-1,\ 0 .
\end{equation}
By the facts just listed, consecutive numbers differ by exactly one. Record the differences as a word of signs, $+$ for
a step up and $-$ for a step down, and mark the step from $\sigma_a$ to $\tau_b-1$ (from $-1$ if $\sigma$ is empty, to
$0$ if $\tau$ is empty) as the \emph{head}. The result is the \emph{tape}: a word of $a+b+1$ letters with one marked
letter, written with the head in brackets. We call the numbers \eqref{eq:path} the \emph{gap numbers}, for the
reason given in the next paragraph. The \emph{root} is the tape of the stack state after crossing $1$: the gap numbers
are $-1,0,-1,0$ and the tape is $+[-]+$.

The tape has a direct picture (Figure~\ref{fig:cut}). The cut meets $a$ open arches below the river, then the river,
then $b$ open arches above. Each of these $a+b+1$ objects carries one letter, read from bottom to top, and the river
carries the head. The numbers \eqref{eq:path} sit in the gaps between the objects: the gap below everything has number
$-1$, crossing an object changes the number by its letter, and the gap above everything has number $0$. The gap just
above an open lower arch carries the label of the crossing where that arch started; the gap just below an open upper
arch carries that label minus one.

The tape determines the two stacks: the gap numbers are the partial sums of the letters, starting from $-1$, and the
head says where $\sigma$ ends and $\tau$ begins. The tape does not always determine $g(j)$: if the head letter is $-$,
then $g(j)=\sigma_a$; if it is $+$, then $g(j)$ can be $\sigma_a$ or $\sigma_a+2$. We call two stack states that
differ only in $g(j)$ \emph{twins}. Twins lead to the same tapes after the next crossing, so nothing is lost by
forgetting $g(j)$.

\subsection{The moves}

Each of the crossings $2,\ldots,2n-1$ becomes one \emph{move} of the tape; the first crossing creates the first tape
and the last one closes the road. Table~\ref{tab:moves} lists the six kinds: two
head moves, three splits, which replace the head letter by three letters, and one merge, which replaces the head and
its two neighbors by one letter. A move reads and changes only the head letter and its two neighbors.

\begin{table}[htb]
\centering
\small
\renewcommand{\arraystretch}{1.2}
\begin{tabular}{llll}
\toprule
move & on the tape & condition & arches at the crossing \\
\midrule
$\Lm$ (head left) & $x[s]y\ \mapsto\ [x]sy$ & head has a left neighbor & lower closes, upper opens\\
$\Rm$ (head right) & $x[s]y\ \mapsto\ xs[y]$ & head has a right neighbor & upper closes, lower opens\\
$\Sone$ (split) & $[+]\ \mapsto\ +[-]+$ & head is $+$ & both open\\
$\Stwo$ (split) & $[-]\ \mapsto\ +[-]-$ & head is $-$ & both open; label up\\
$\Sthree$ (split) & $[-]\ \mapsto\ -[-]+$ & head is $-$ & both open; label down\\
$\Mm$ (merge) & $x[+]y\ \mapsto\ [x+y+1]$ & two neighbors, not $+[+]+$ & both close\\
\bottomrule
\end{tabular}
\caption{The moves of the tape machine. Here $x$, $s$, $y$ stand for letters, and the rest of the tape is
unchanged. In the merge the letters count as $\pm1$: $-[+]-$ becomes $[-]$, and
$-[+]+$ and $+[+]-$ become $[+]$. The replacement at $+[+]+$ would give a step of $3$ and does not exist. At most four
moves apply to any tape.}
\label{tab:moves}
\end{table}

The table follows from the stack rules by cases. If the head letter is $-$, both innermost labels equal $g(j)$. A step
up must open an upper arch (an upper closing would need label $g(j)+2$); below it either opens an arch, which gives
$\Stwo$, or closes one, which moves the head one letter to the left. A step down is the mirror image: $\Sthree$ or a
move of the head to the right. If the head letter is $+$, the innermost upper label is two more than the innermost
lower label $\sigma_a$, and the next label must be $\sigma_a+1$. Opening on both sides gives $\Sone$; closing below and
opening above moves the head left; opening below and closing above moves it right; closing on both sides deletes the
two innermost entries, and three consecutive steps of \eqref{eq:path} become one. Their sum is $\pm1$ except when
all three are $+$; then the sum is $3$, and the new state would need the next label to go down and up at once, so it
has no continuation. This is the excluded merge at $+[+]+$. Finally, an empty stack corresponds exactly to a head at
an end of the tape, which explains the conditions on the head moves and the merge. Figure~\ref{fig:moves} shows the
moves at two tapes of the running example.

\begin{figure}[htb]
\centering
\begin{tikzpicture}[>=Stealth]
  % panel (a): head on -
  \tp{(0,0)}{+,-,+,-,+}{4}
  \node[font=\small,anchor=east] at (-0.1,0) {(a)};
  \tp{(3.0,1.5)}{+,-,+,-,+}{3}
  \tp{(3.0,0.5)}{+,-,+,-,+}{5}
  \tp{(3.0,-0.5)}{+,-,+,+,-,-,+}{5}
  \tp{(3.0,-1.5)}{+,-,+,-,-,+,+}{5}
  \foreach \yy/\lab in {1.5/{$\Lm$},0.5/{$\Rm$},-0.5/{$\Stwo$},-1.5/{$\Sthree$}} {
    \draw[->,line width=0.5pt] (1.95,0)--(2.85,\yy) node[pos=0.62,fill=white,inner sep=1pt,font=\scriptsize] {\lab};}
  % panel (b): head on +
  \begin{scope}[xshift=8.0cm]
  \tp{(0,0)}{+,-,+,-,+}{3}
  \node[font=\small,anchor=east] at (-0.1,0) {(b)};
  \tp{(3.0,1.5)}{+,-,+,-,+}{2}
  \tp{(3.0,0.5)}{+,-,+,-,+}{4}
  \tp{(3.0,-0.5)}{+,-,+,-,+,-,+}{4}
  \tp{(3.0,-1.5)}{+,-,+}{2}
  \foreach \yy/\lab in {1.5/{$\Lm$},0.5/{$\Rm$},-0.5/{$\Sone$},-1.5/{$\Mm$}} {
    \draw[->,line width=0.5pt] (1.95,0)--(2.85,\yy) node[pos=0.62,fill=white,inner sep=1pt,font=\scriptsize] {\lab};}
  \end{scope}
\end{tikzpicture}
\caption{The moves at two tapes. (a) The tape $+-+[-]+$ of Figure~\ref{fig:cut}, with the head on a $-$: two head
moves and two splits. (b) The tape $+-[+]-+$ reached after crossing $4$, with the head on a $+$: two head moves, one
split and the merge $-[+]-\mapsto[-]$, which returns to the root. The shaded box is the head.}
\label{fig:moves}
\end{figure}

A \emph{walk} is a sequence of tapes in which each tape is obtained from the previous one by a move; its length is
the number of moves. A \emph{closed walk} is a walk that starts and ends at the root.

\begin{theorem}[tape machine]\label{thm:tape}
For every $n\ge1$, reading the crossings $2,\ldots,2n-1$ of a meander as moves is a bijection between the meanders of
order $n$ and the walks of $2n-2$ moves of Table~\ref{tab:moves} that start at the root $+[-]+$ and end at the root.
\end{theorem}

\begin{proof}[Idea of the proof]
\emph{Labels.} Deleting crossing $2n$ turns the road into a path, and Lemma~\ref{lem:labels} gives labels that change
by $\pm1$ along arches and between neighbors along the river.

\emph{Stacks and differences.} At each cut the open arches form two stacks of labels, and the numbers
\eqref{eq:path} form a path with steps $\pm1$; its steps and the marked step are the tape. The tape recovers both
stacks.

\emph{One crossing, one move.} Revealing the next crossing pushes or pops each stack, and the case analysis above
shows that the possible updates are exactly the moves of Table~\ref{tab:moves}, that different updates give different
tapes, and that twin stack states lead to the same tapes. Hence a walk of tapes starting at the root lifts in exactly
one way to a walk of stack states, and it ends at the first stack state if and only if the tape walk ends at the
root (the root has head letter $-$, so its label is determined). The first crossing creates the root and the last
crossing closes the two remaining arches, so a meander of order $n$ gives a closed walk of $2n-2$ moves.

\emph{Back to a meander.} Conversely, replay a closed walk from the root, remembering at which crossing each arch was
opened. Every closing joins the current crossing to the most recently opened arch on its side, so the arches do not
cross; at the end one arch remains on each side and both are joined to a final crossing $2n$. The labels obey
\eqref{eq:arch} and the neighbor rule by construction, so by the converse half of Lemma~\ref{lem:labels} the arches
form one closed road. The two constructions undo each other.
\end{proof}

The same proof works on a strip: closed walks whose gap numbers stay in a fixed range of $K$ consecutive integers
count the meanders whose labels stay in the same range extended by one value at the top, a fixed band of $K+1$
integers. For example, gap numbers from $-1$ to $1$ correspond to labels from $-1$ to $2$. Strips are used for the
lower bound (Section~\ref{sec:flows}).

\paragraph{The excluded merge and the one-road property.}
The machine never closes a loop too early. Such a step would be a crossing that closes, on both sides, two arches
belonging to the same piece of road. No allowed move does this. Every tape reached from the root can be completed to a closed walk
(Lemma~\ref{lem:sweep} below), every closed walk gives a single road (the converse half of
Lemma~\ref{lem:labels}), and a loop closed early would survive in the completion as a second component. At a head
$-$ the machine has no merge at all: there the two innermost arches carry the same label, and the label of the new
crossing would have to be one more and one less than that label at once. The merge at $+[+]+$ is missing for a
different reason, label consistency: after it the arch below would force the next label down and the arch above would
force it up, so no continuation exists. The local rules know nothing about connectivity; it is the planar argument of
Lemma~\ref{lem:labels} that turns the label rules into the global condition that the road is one closed curve.
% Author's check, not in the sources (script written for this paper): over all histories of at most 11 moves, at every
% tape with head letter + the two innermost open arches belong to different pieces of road, including at + [+] +.
% UNVERIFIED as a general statement; not used anywhere.

\subsection{Worked examples}

\emph{Order $2$.} The root has four moves: $\Lm$, $\Rm$, $\Stwo$, $\Sthree$. After a split the head is on a $-$,
where no merge applies, so only the head moves can be undone in one more move. The closed walks of $2$ moves are
therefore $\Lm\Rm$ and $\Rm\Lm$:
\[ +[-]+ \xrightarrow{\ \Lm\ } [+]-+ \xrightarrow{\ \Rm\ } +[-]+, \qquad
   +[-]+ \xrightarrow{\ \Rm\ } +-[+] \xrightarrow{\ \Lm\ } +[-]+ . \]
The first gives the lower arches $1$--$2$, $3$--$4$ and the upper arches $1$--$4$, $2$--$3$; the second is its mirror
image. These are the $M_2=2$ meanders.

\emph{Order $3$.} The running example is the closed walk $\Rm\,\Sone\,\Lm\,\Mm$. Table~\ref{tab:example} follows it
crossing by crossing.

\begin{table}[htb]
\centering
\small
\setlength{\tabcolsep}{4pt}
\begin{tabular}{cclcccl}
\toprule
crossing & label & arches at this crossing & $\sigma$ & $\tau$ & tape after it & move \\
\midrule
$1$ & $0$ & opens lower $1$--$6$, upper $1$--$2$ & $(0)$ & $(0)$ & $+[-]+$ & (start)\\
$2$ & $-1$ & closes upper $1$--$2$, opens lower $2$--$5$ & $(0,-1)$ & $()$ & $+-[+]$ & $\Rm$\\
$3$ & $0$ & opens lower $3$--$4$, upper $3$--$6$ & $(0,-1,0)$ & $(0)$ & $+-+[-]+$ & $\Sone$\\
$4$ & $1$ & closes lower $3$--$4$, opens upper $4$--$5$ & $(0,-1)$ & $(0,1)$ & $+-[+]-+$ & $\Lm$\\
$5$ & $0$ & closes lower $2$--$5$ and upper $4$--$5$ & $(0)$ & $(0)$ & $+[-]+$ & $\Mm$\\
$6$ & & closes lower $1$--$6$ and upper $3$--$6$ & & & & (end)\\
\bottomrule
\end{tabular}
\caption{The running example as a walk of the tape machine. The merge at crossing $5$ is $-[+]-\mapsto[-]$.}
\label{tab:example}
\end{table}

The other seven meanders of order $3$ are the walks $\Lm\Rm\Lm\Rm$, $\Lm\Rm\Rm\Lm$, $\Rm\Lm\Lm\Rm$, $\Rm\Lm\Rm\Lm$,
$\Lm\,\Sone\,\Rm\,\Mm$, $\Stwo\,\Lm\,\Mm\,\Rm$ and $\Sthree\,\Rm\,\Mm\,\Lm$. Counting closed walks by computer gives
\[ 1,\ 2,\ 8,\ 42,\ 262,\ 1828,\ 13820,\ 110954,\ \ldots \]
for $n=1,2,3,\ldots$, the meander numbers; the counts were compared with the exact table up to $n=15$. A
correct count for small $n$ is a check of the implementation, not a proof; the proof is Theorem~\ref{thm:tape}.

\subsection{Basic facts}

Let $\cls$ be the class of tapes with at least three letters whose first and last letters are $+$ and whose letters
sum to $1$ (counting $+$ as $+1$ and $-$ as $-1$). The root is in $\cls$, and every move maps $\cls$ to itself: a split
replaces a letter by three with the same sum, a merge replaces three by one with the same sum, and neither changes an
end letter $+$ into $-$. A merge never acts on a tape of three letters, because the only such tape in $\cls$ is $+-+$,
whose middle letter is $-$.
Let $a$ and $b$ be the numbers of letters to the left and to the right of the head. The moves change $(a,b)$ by
\begin{equation}\label{eq:drift}
  \Lm:(-1,+1),\qquad \Rm:(+1,-1),\qquad \text{splits}:(+1,+1),\qquad \Mm:(-1,-1).
\end{equation}
So $a$ changes parity at every move, and a closed walk has an even number of moves.

\begin{lemma}[linear completion]\label{lem:sweep}
Every tape of $\cls$ with $N$ letters can be reached from the root in at most $3N$ moves, and the root can be reached
from it in at most $3N$ moves, through tapes with at most $N$ letters.
\end{lemma}

\begin{proof}
A merge deletes the head letter $+$ and one neighboring $-$, and leaves the other neighbor in place. A split read
backward deletes the head letter $-$ and one neighboring $+$. Put the head on the second letter and sweep to the right,
keeping the inner letters before the head all equal. When the head and the inner letter before it differ, delete this
pair, by a merge for the return and by a backward split for the build, after at most one head move; the head moves
back by at most two places. Otherwise move right. When the head reaches the last letter, all inner letters are equal
and sum to $-1$, so the tape is $+-+$. There are $(N-3)/2$ deletions; counting the moves gives at most $3N$. Read
backward, the build is a legal walk from the root, because head moves are reversible and a backward split read forward
is a split.
\end{proof}

Closed walks can be concatenated: one closed walk followed by another is again a closed walk. Write $c_t$ for the
number of closed walks of $t$ moves, so that $c_{2n-2}=M_n$. Then $c_tc_{t'}\le c_{t+t'}$, and since at most four moves
apply at any tape, $c_t\le 4^t$. By Fekete's lemma on supermultiplicative sequences
\begin{equation}\label{eq:fekete}
  \mu=\lim_{n\to\infty}c_{2n}^{1/n}=\sup_n c_{2n}^{1/n},\qquad\text{so}\qquad c_t\le\mu^{t/2}\ \text{for every even }t.
\end{equation}
This is also why squares appear in both bounds: one unit of $n$ is two moves.
